% bare_jrnl.tex
% V1.4b
% 2015/08/26
% by Michael Shell
% see http://www.michaelshell.org/
% for current contact information.
%
% This is a skeleton file demonstrating the use of IEEEtran.cls
% (requires IEEEtran.cls version 1.8b or later) with an IEEE
% journal paper.
%
% Support sites:
% http://www.michaelshell.org/tex/ieeetran/
% http://www.ctan.org/pkg/ieeetran
% and
% http://www.ieee.org/

%%*************************************************************************
%% Legal Notice:
%% This code is offered as-is without any warranty either expressed or
%% implied; without even the implied warranty of MERCHANTABILITY or
%% FITNESS FOR A PARTICULAR PURPOSE!
%% User assumes all risk.
%% In no event shall the IEEE or any contributor to this code be liable for
%% any damages or losses, including, but not limited to, incidental,
%% consequential, or any other damages, resulting from the use or misuse
%% of any information contained here.
%%
%% All comments are the opinions of their respective authors and are not
%% necessarily endorsed by the IEEE.
%%
%% This work is distributed under the LaTeX Project Public License (LPPL)
%% ( http://www.latex-project.org/ ) version 1.3, and may be freely used,
%% copied, modified and distributed. A copy of the LPPL, version 1.3, is
%% included in the base LaTeX documentation of all distributions of LaTeX
%% released 2003/12/01 or later.
%% Retain all contribution notices and credits.
%% ** Modified files should be clearly indicated as such, including  **
%% ** renaming them and changing author support contact information. **
%%*************************************************************************


% *** Authors should verify (and, if needed, correct) their LaTeX system  ***
% *** with the testflow diagnostic prior to trusting their LaTeX platform  ***
% *** with production work. The IEEE's font choices and paper sizes can    ***
% *** trigger bugs that do not appear when using other class files.        ***
% The testflow support page is at:
% http://www.michaelshell.org/tex/testflow/


\documentclass[journal]{IEEEtran}
% Some very useful LaTeX packages include:
% (uncomment the ones you want to load)


% *** MATH PACKAGES ***
%
\usepackage{amsmath}
% A popular package from the American Mathematical Society that provides
% many useful and powerful commands for dealing with mathematics.
%
% Note that the amsmath package sets \interdisplaylinepenalty to 10000
% thus preventing page breaks from occurring within multiline equations. Use:
%\interdisplaylinepenalty=2500
% after loading amsmath to restore such page breaks as IEEEtran.cls normally
% does. amsmath.sty is already installed on most LaTeX systems. The latest
% version and documentation can be obtained at:
% http://www.ctan.org/pkg/amsmath


% *** HYPERLINKS ***
%
\usepackage{hyperref}
% Provides \href{url}{text} for clickable links in PDF output.


% correct bad hyphenation here
\hyphenation{op-tical net-works semi-conduc-tor}


\begin{document}

% paper title
% Titles are generally capitalized except for words such as a, an, and, as,
% at, but, by, for, in, nor, of, on, or, the, to and up, which are usually
% not capitalized unless they are the first or last word of the title.
% Linebreaks \\ can be used within to get better formatting as desired.
% Do not put math or special symbols in the title.
\title{Two-Stage Visual-Inertial SLAM with Sliding-Window Filtering and
Global Bundle Adjustment for the Hilti $\times$ Trimble SLAM Challenge 2026}

% author names and IEEE memberships
% note positions of commas and nonbreaking spaces ( ~ ) LaTeX will not break
% a structure at a ~ so this keeps an author's name from being broken across
% two lines.
% use \thanks{} to gain access to the first footnote area
% a separate \thanks must be used for each paragraph as LaTeX2e's \thanks
% was not built to handle multiple paragraphs
\author{Arda~C.~Demirta\c{s}$^{*1}$,
        Berkay~\c{S}ekeroglu$^{*1}$,
        and~Atakan~Topaloglu$^{*1}$%
\thanks{$^{*}$All authors contributed equally to this submission.}%
\thanks{$^{1}$The authors are with the Department of Information Technology
  and Electrical Engineering (D-ITET), ETH Z\"urich, 8092 Z\"urich,
  Switzerland. \texttt{\{ademirtas, bsekeroglu, atopaloglu\}@ethz.ch}}%
\thanks{Manuscript submitted to the Hilti $\times$ Trimble
  SLAM Challenge 2026.}}

% make the title area
\maketitle

% As a general rule, do not put math, special symbols or citations
% in the abstract or keywords.
\begin{abstract}
This report summarizes our submission to the Hilti $\times$ Trimble SLAM Challenge 2026,
covering both the SLAM and Localization tasks.
We developed a two-stage visual-inertial SLAM pipeline: first, a causal
multi-camera MSCKF-based estimator produces an initial trajectory and feature
tracks; second, an offline visual-inertial bundle adjustment refines all
trajectories using the full sequence.
During development, the main bottleneck was not only drift, but reliable
initialization and sustained feature tracking in dynamic-start, low-texture,
and low-light sequences.
Our final SLAM submission relies on dynamic initialization, reduced
playback rate for stable tracking, and per-sequence selection among a small
set of robustness-oriented filter configurations.
For the Localization task, we register the SLAM trajectories to the floorplan
coordinate frame using a single cam0 anchor pose provided per sequence,
applying an SE(2) transform derived from the IMU-to-camera extrinsic calibration
and the ground-truth initialisation pose.
\end{abstract}

% Note that keywords are not normally used for peerreview papers.
\begin{IEEEkeywords}
Visual-inertial odometry, SLAM, MSCKF, bundle adjustment, IMU preintegration, OpenVINS, COLMAP,
floorplan localization, map alignment, Hilti SLAM Challenge.
\end{IEEEkeywords}


% For peer review papers, you can put extra information on the cover
% page as needed:
% \ifCLASSOPTIONpeerreview
% \begin{center} \bfseries EDICS Category: 3-BBND \end{center}
% \fi
%
% For peerreview papers, this IEEEtran command inserts a page break and
% creates the second title. It will be ignored for other modes.
\IEEEpeerreviewmaketitle


\section{Introduction}
% The very first letter is a 2 line initial drop letter followed
% by the rest of the first word in caps.
%
% form to use if the first word consists of a single letter:
% \IEEEPARstart{A}{demo} file is ....
%
% form to use if you need the single drop letter followed by
% normal text (unknown if ever used by IEEE):
% \IEEEPARstart{A}{}
%
% Some journals put the first two words in caps:
% \IEEEPARstart{T}{his demo} file is ....
%
% Here we have the typical use of a "T" for an initial drop letter
% and "HIS" in caps to complete the first word.
\IEEEPARstart{T}{his} report describes our final approach for the Hilti
$\times$ Trimble SLAM Challenge 2026 trajectory-estimation task.
Since the SLAM task evaluates trajectories in any preferred reference frame,
we focused on producing geometrically consistent trajectories with high
temporal coverage and low drift.

The dataset exposed two dominant issues.
First, several sequences were difficult to initialize because they began with
motion, low visual texture, or rapid viewpoint changes.
Second, even when initialization succeeded, feature-poor regions caused drift
that a causal sliding-window estimator could not fully recover.

Our final pipeline addresses these issues with a two-stage design.
We use OpenVINS \cite{geneva2020openvins} as the front-end estimator because
it provides robust MSCKF-based visual-inertial filtering and exposes the
feature tracks required for later refinement.
We then apply an offline visual-inertial bundle adjustment stage to all
trajectories to reduce residual drift, using the
\href{https://github.com/B1ueber2y/colmap/tree/features/imu}{COLMAP visual-inertial fork}
\cite{schoenberger2016sfm}.

The final system consists of five steps:
\begin{enumerate}
    \item Run an MSCKF visual-inertial filter with KLT tracking across both
          fisheye cameras and IMU propagation at 1000\,Hz.
    \item Generate multiple candidate trajectories using robustness-oriented
          filter configurations.
    \item Select the best candidate trajectory for each sequence.
    \item Reconstruct feature tracks from the filter output.
    \item Refine all trajectories using offline visual-inertial bundle adjustment.
\end{enumerate}



\section{Filter Initialization and Configuration}

\subsection{Initialization Failure Modes}

Initialization was the primary failure mode across the 30 sequences.
The static initializer requires feature tracks to survive for at least half
the initialization window in the older half of the sliding window before
attempting state recovery.
Fast motion or low visual texture breaks KLT tracks before they age
sufficiently, blocking initialization indefinitely.

A representative case is \texttt{floor\_EG/2025-12-02/run\_2}.
The static initializer waits for a still-to-moving jerk event that does not
occur until approximately 59\,s into the recording, when the operator
adjusts their grip.
Switching to dynamic initialization with a 1.0\,s window allowed the filter
to trigger on the early camera-shake motion, reducing the initialization
delay to 5.6\,s.

\subsection{Dynamic Initialization}

Dynamic initialization triggers on generally excited angular velocity rather
than requiring a static-to-moving transition.
This was the single most effective intervention, converting several delayed
or failed runs into complete trajectories.
It is particularly important for sequences that begin while the operator is
already walking.

\subsection{Initialization Window Duration}

A shorter initialization window (1.0\,s versus the 3.0\,s default) reduces
the track-survival requirement at a minor cost to inertial alignment quality.
On dynamic-start and low-texture sequences this trade-off was clearly
beneficial.

\subsection{Reduced Playback Rate}

At full playback speed, the CPU could not consistently process every image
frame at the nominal 30\,Hz camera rate, reaching only 16--23 frames per
second.
This forced the optical-flow tracker to operate on larger apparent
inter-frame displacements, shortening feature tracks and degrading both
initialization and sustained estimation.
Processing at 0.5$\times$ playback speed restored near-nominal tracking
performance and increased pose output density from approximately 16--23\,Hz
to approximately 29\,Hz.

\subsection{Robustness-Oriented Filter Settings}

We designed five filter variants around the baseline configuration.
The useful changes were: increasing the number of tracked features and
persistent SLAM landmarks; extending the clone horizon; applying CLAHE
contrast normalization; inflating IMU process noise; and enabling
zero-velocity updates (ZUPT).

Inflating the IMU process noise covariance down-weights the inertial
contribution to the EKF update, allowing visual measurements to dominate
when inertial biases are poorly observable (for example, slow walking in
low-contrast basement corridors).
The trade-off is increased susceptibility to IMU integration drift under
fast motion.

Increasing tracked features to 800 and SLAM landmarks to 100 widens the
reprojection constraint set.
Extending the clone horizon to 15 allows the filter to accumulate longer-baseline
observations, improving triangulation depth in texture-sparse regions.

ZUPT constrains velocity and rotation rate to zero during detected stationary
intervals, strongly anchoring bias estimation.
In practice it had limited impact because most sequences involve continuous
walking with few clear stationary periods.


\section{Per-Sequence Selection}

No single filter configuration was best for all sequences.
We generated multiple candidate trajectories per run and selected the final
one using three criteria.

For the five released ground-truth sequences, selection was based on ATE after
SE(3) alignment.
For the remaining sequences, we used trajectory plausibility and
multi-configuration consensus.
Specifically, we checked whether the horizontal extent was consistent with the
expected floor geometry and whether vertical drift remained below approximately
2\,m (elevated vertical drift reliably indicates unconstrained IMU bias
divergence in monocular mode).
When three or more independent configurations converged on the same trajectory
shape, that shape was accepted as correct.

Table~\ref{tab:config_distribution} summarizes the final selected
configuration types.

% An example of a floating table. Note that, for IEEE style tables, the
% \caption command should come BEFORE the table and, given that table
% captions serve much like titles, are usually capitalized except for words
% such as a, an, and, as, at, but, by, for, in, nor, of, on, or, the, to
% and up, which are usually not capitalized unless they are the first or
% last word of the caption. Table text will default to \footnotesize as
% the IEEE normally uses this smaller font for tables.
% The \label must come after \caption as always.
\begin{table}[!t]
\renewcommand{\arraystretch}{1.2}
\caption{Final Submission Configuration Types}
\label{tab:config_distribution}
\centering
\begin{tabular}{l c}
\hline
Configuration type & Runs \\
\hline
Dynamic init, 1.0\,s window, 0.5$\times$ playback & 13 \\
Dynamic init, 3.0\,s window, 0.5$\times$ playback & 10 \\
Baseline, full playback speed & 3 \\
Inflated IMU process noise, 0.5$\times$ playback & 2 \\
High-feature, extended clone horizon, 0.5$\times$ playback & 2 \\
\hline
\end{tabular}
\end{table}

Dynamic initialization at 0.5$\times$ playback accounts for 77\% of the
final submission.


\section{Offline Visual-Inertial Bundle Adjustment}

We applied an offline visual-inertial bundle adjustment to all runs to
refine the filter output.
The optimization is seeded directly from the filter: the feature publisher
outputs, at each camera frame, the 3D position, pixel coordinates, and
stable track ID of every active KLT-tracked feature.
Accumulating these messages reconstructs the full track history without
additional feature extraction or descriptor matching.
A typical run yields approximately 20k--50k triangulated 3D points and
800k 2D reprojection observations.

The bundle adjustment minimizes the joint cost
\begin{equation}
\min_{\mathbf{T},\mathbf{p},\mathbf{b}}
\sum_{(i,j)} \rho\!\left(\left\|\mathbf{r}^{ij}_{\mathrm{reproj}}\right\|^2_{\Sigma_c}\right)
+ \sum_k \left\|\mathbf{r}^{k}_{\mathrm{IMU}}\right\|^2_{\Sigma_{\mathrm{IMU}}},
\label{eq:viba}
\end{equation}
where $\mathbf{T}$ are camera poses, $\mathbf{p}$ are landmark positions,
$\mathbf{b}$ denotes inertial states and biases,
$\mathbf{r}^{ij}_{\mathrm{reproj}}$ are per-observation reprojection
residuals, $\mathbf{r}^{k}_{\mathrm{IMU}}$ are IMU preintegration residuals
computed using the on-manifold analytical formulation of Forster et
al.~\cite{forster2017preintegration}, and $\rho(\cdot)$ is a Huber loss.
The optimization variables are refined jointly without fixing the filter
poses.

The principal limitation is that KLT tracks carry no feature descriptors.
When a track is lost due to occlusion or fast motion, the physical point
re-appears with a new track ID, breaking the observation chain.
This weakens the global reprojection term, particularly in environments with
heavy occlusion or repetitive structure.

Table~\ref{tab:viba_result} summarises the aggregate VI-BA improvement
over the filter baseline on the five ground-truth sequences.

\begin{table}[!t]
\renewcommand{\arraystretch}{1.2}
\caption{Effect of Offline VI-BA (averaged over 5 GT sequences)}
\label{tab:viba_result}
\centering
\begin{tabular}{l c}
\hline
Method & ATE RMSE (m) \\
\hline
Filter baseline (mean, 5 GT runs) & 0.327 \\
After VI-BA (mean, 5 GT runs) & \textbf{0.262} \\
Mean score improvement & \textbf{+2.4\,pts} \\
\hline
\end{tabular}
\end{table}


\section{Results}

Table~\ref{tab:gt_results} reports per-run performance on the five released
ground-truth sequences after SE(3) alignment.

\begin{table}[!t]
\renewcommand{\arraystretch}{1.2}
\caption{Per-Run Results on the Five Ground-Truth Sequences}
\label{tab:gt_results}
\centering
\begin{tabular}{l c c c c}
\hline
Sequence & \multicolumn{2}{c}{ATE RMSE (m)} & \multicolumn{2}{c}{Score} \\
         & Filter & VI-BA & Filter & VI-BA \\
\hline
floor\_1 / 2025-05-05 / run\_1   & 0.300 & \textbf{0.233} & 85.7 & \textbf{88.3} \\
floor\_2 / 2025-05-05 / run\_1   & 0.446 & \textbf{0.390} & 81.4 & \textbf{83.3} \\
floor\_2 / 2025-10-28 / run\_1   & 0.248 & \textbf{0.184} & 89.7 & \textbf{92.5} \\
floor\_2 / 2025-10-28 / run\_2   & 0.262 & \textbf{0.192} & 89.5 & \textbf{91.9} \\
floor\_UG1 / 2025-10-16 / run\_1 & 0.378 & \textbf{0.311} & 82.5 & \textbf{85.1} \\
\hline
\textbf{Mean}                    & 0.327 & \textbf{0.262} & 85.8 & \textbf{88.2} \\
\hline
\end{tabular}
\end{table}

The weakest result is \texttt{floor\_2/2025-05-05/run\_1}, which contains a
long feature-poor corridor where a monocular visual-inertial filter without
loop closure accumulates translational drift.


\section{Localization Task}

Each SLAM trajectory is expressed in an arbitrary global frame; we register
it to the floorplan map using the single ground-truth cam0 anchor pose
$\mathbf{T}^{\text{map}}_{\text{cam}_0,\text{gt}}$ provided per sequence.

SLAM poses are output in the IMU frame, so the anchor SLAM pose is first
converted to cam0 using the Kalibr extrinsic $\mathbf{T}^{\text{imu}}_{\text{cam}_0}$:
\begin{equation}
  \mathbf{T}^{\text{global}}_{\text{cam}_0,k}
    = \mathbf{T}^{\text{global}}_{\text{imu},k}\,\mathbf{T}^{\text{imu}}_{\text{cam}_0}.
\end{equation}
The global-to-map transform follows as
\begin{equation}
  \mathbf{T}^{\text{map}}_{\text{global}}
    = \mathbf{T}^{\text{map}}_{\text{cam}_0,\text{gt}}\,
      \bigl(\mathbf{T}^{\text{global}}_{\text{cam}_0,k}\bigr)^{-1}.
\end{equation}
Because the SLAM estimator maintains a gravity-aligned frame, the residual
pitch and roll of $\mathbf{T}^{\text{map}}_{\text{global}}$ are noise.
We zero them, retaining only the yaw $\psi$ and translation $\mathbf{t}$,
and apply the resulting planar transform to every pose in the trajectory:
\begin{equation}
  \mathbf{T}^{\text{map}}_{\text{cam}_0,i}
    = \begin{bmatrix} \mathbf{R}_z(\psi) & \mathbf{t} \\ \mathbf{0}^\top & 1 \end{bmatrix}
      \mathbf{T}^{\text{global}}_{\text{imu},i}\,\mathbf{T}^{\text{imu}}_{\text{cam}_0}.
\end{equation}
For sequences where the filter initializes after the anchor epoch, we
decouple the rotation and translation estimates: the yaw is derived from
the first filter-initialized pose, while the translation is set so that
the anchor-epoch SLAM position maps exactly to the ground-truth map position.


\section{Discussion}

Robust initialization mattered more than back-end complexity.
Dynamic initialization was the most impactful change, converting delayed or
failed runs into complete trajectories.
The reduced playback rate was equally important in practice: processing frames
faster than the optical-flow tracker can reliably handle degrades track
quality and ultimately causes filter divergence, even though the underlying
data is identical.

The offline bundle-adjustment stage confirmed that full-trajectory
optimization recovers drift that is irrecoverable in a sliding-window filter.
The main limitation is the reliance on KLT tracks without descriptors.
Replacing the KLT front-end with descriptor-based local features and
re-identification across track breaks (for example, ALIKED with LightGlue
matching) would yield longer, more constraining observation chains and
likely improve the back-end substantially, particularly in the occluded
basement environments.


\section{Conclusion}

We presented a two-stage visual-inertial SLAM pipeline combining an MSCKF
sliding-window filter with offline global visual-inertial bundle adjustment.
The most effective practical decisions were dynamic initialization, reduced
playback speed for stable feature tracking, and per-sequence selection among
a small number of targeted filter variants.
For the Localization task, SLAM trajectories are registered to the floorplan
frame via a gravity-consistent planar transform derived from the IMU-to-camera
extrinsic calibration and a single ground-truth anchor pose per sequence.


% use section* for acknowledgment
\section*{Acknowledgment}

The authors thank their supervisors Zador Pataki
(\texttt{zador.pataki@inf.ethz.ch}), Xudong Jiang
(\texttt{xudong.jiang@inf.ethz.ch}), and Shaohui Liu
(\texttt{shaohui.liu@inf.ethz.ch}) of the Department of Computer Science,
ETH Z\"urich, and Paul-Edouard Sarlin (\texttt{psarlin@google.com}) of
Google, for their guidance throughout this project.
We also thank the organizers of the Hilti $\times$ Trimble SLAM Challenge 2026.


% trigger a \newpage just before the given reference
% number - used to balance the columns on the last page
% adjust value as needed - may need to be readjusted if
% the document is modified later
%\IEEEtriggeratref{8}
% The "triggered" command can be changed if desired:
%\IEEEtriggercmd{\enlargethispage{-5in}}

% references section

% can use a bibliography generated by BibTeX as a .bbl file
% BibTeX documentation can be easily obtained at:
% http://mirror.ctan.org/biblio/bibtex/contrib/doc/
% The IEEEtran BibTeX style support page is at:
% http://www.michaelshell.org/tex/ieeetran/bibtex/
%\bibliographystyle{IEEEtran}
% argument is your BibTeX string definitions and bibliography database(s)
%\bibliography{IEEEabrv,../bib/paper}
%
% <OR> manually copy in the resultant .bbl file
% set second argument of \begin to the number of references
% (used to reserve space for the reference number labels box)
\begin{thebibliography}{1}

\bibitem{geneva2020openvins}
P.~Geneva, K.~Eckenhoff, W.~Lee, Y.~Yang, and G.~Huang,
``OpenVINS: A research platform for visual-inertial estimation,''
in \emph{Proc. IEEE Int. Conf. Robot. Autom. Workshop on Open Source
Systems in Robotics}, 2020. \texttt{https://github.com/rpng/open\_vins}

\bibitem{schoenberger2016sfm}
J.~L. Sch\"onberger and J.-M. Frahm,
``Structure-from-motion revisited,''
in \emph{Proc. IEEE Conf. Comput. Vis. Pattern Recognit.}, 2016,
pp.~4104--4113.

\bibitem{forster2017preintegration}
C.~Forster, L.~Carlone, F.~Dellaert, and D.~Scaramuzza,
``On-manifold preintegration for real-time visual-inertial odometry,''
\emph{IEEE Trans. Robot.}, vol.~33, no.~1, pp.~1--21, 2017.

\end{thebibliography}

% biography section
%
% If you have an EPS/PDF photo (graphicx package needed) extra braces are
% needed around the contents of the optional argument to biography to prevent
% the LaTeX parser from getting confused when it sees the complicated
% \includegraphics command within an optional argument. (You can create
% your own custom biographies by commenting out the above and using the
% unindented \begin{biography} form instead to achieve this.)
%
% \begin{IEEEbiography}[{\includegraphics[width=1in,height=1.25in,clip,keepaspectratio]{mshell}}]{Michael Shell}
% Biography text here.
% \end{IEEEbiography}

% if you will not have a photo at all:
% \begin{IEEEbiographynoskip}{Michael Shell}
% Biography text here.
% \end{IEEEbiographynoskip}

% insert where needed to balance the two columns on the last page with
% biographies
%\newpage

% You can push biographies down or up by placing
% a \vfill before or after them. The appropriate
% use of \vfill depends on what kind of text is
% on the last page and whether or not the columns
% are being equalized.

%\vfill

% Can be used to pull up biographies so that the bottom of the last one
% is flush with the other column.
%\enlargethispage{-5in}

\end{document}
